\scanpage{124}
% archive sheet number: 123
% source lines 1--18
le foncteur compatible aux $\otimes$ associatifs unitaires qu'on obtient de $C$ dans $\mathrm{Rep}_k(G)$, et voir en particulier sous quelles conditions il est une équivalence de catégories: \deleted{Alors}\inserted{\handwritten{il \unsure{f. et s.|f. et p.} que $A$ soit commutatif! ``\unsure{Géométriquement|biométriquement}''}} Lorsqu'il en est ainsi,\inserted{\handwritten{\uuline{dans le cas gradué anticomm.}}} la situation est décrite entièrement par la donnée de $G$, de la graduation sur \deleted{\unsure{son}} le foncteur oubli sur \deleted{$C$}\inserted{$R_k(G)$}, décrite par un homomorphisme de monoïdes $\underline{G}_m\handwritten{\to}G$, et enfin la condition que pour deux représentation $M,N$ de $G$, l'isomorphisme de Koszul $M\otimes N\handwritten{\xrightarrow{\sim}}N\otimes M$ sur les vectoriels sous-jacents est compatible avec l'action de $G$. Bien sûr, si $k$ est de car 2, ce n'est pas une restriction, et dans ce cas d'ailleurs Koko ne change rien i.e. tous nos ennuis et distinguos disparaissent. \deleted{Si} Il suffit en tous cas que l'homomorphisme induit $\handwritten{\mu}_2\handwritten{\longrightarrow}G$ soit central;\inserted{\handwritten{et on voit que cela est \unsure{vérifié|vrai} si \unsure{(car $k\ne2$, i.e.)|(car $k\ne2$, c.à.d.)} \unsure{que cela implique} $A^{\unsure{*|\#}}=A^{\unsure{\oplus|0}}$ donc $A^*$ commutatif !}} à fortiori\inserted{\handwritten{il suffit}} que $\underline{G}_m\handwritten{\to}G$ soit central\handwritten{.}\deleted{,i.e.}\inserted{Cela signifie} que la décomposition d'un $M$ en composants \deleted{homogènes}\inserted{pairs et impairs} soit invariante sous $G$ i.e. se fasse dans $C$,resp. qu'il en soit de même pour la décomposition en degrés\inserted{\handwritten{entiers}}. Noter que cette dernière condition s'exprime directement en termes de $C$ (muni de sa structure $\otimes$ et du foncteur \deleted{oubli} gradué\deleted{\unsure{x|s}} $C\handwritten{\to}\mathrm{Grdf}(k)$), \deleted{et on \deleted{espère}\inserted{\handwritten{voit}} que cette condition est suffisante (peut-être nécessaire et suffisante)\inserted{\handwritten{si car. $k\ne2$}} pour qu'on soit dans le cas favorable\inserted{\handwritten{où}} $C\handwritten{\to}\mathrm{Rep}_k(G)$ \unsure{est} une équivalence.}
\margin{\deleted{\handwritten{et on \unsure{voit}}}}

% source lines 19--27
\handwritten{10)} Ce résultat s'appliquera alors, en particulier, à la catégorie des motifs\inserted{semi-simples} (car les \deleted{foncteurs} graduations naturelles sur les foncteurs $H^\bullet$ qu'on arrive à y définir \deleted{sont} proviennent d'une graduation de la catégorie des motifs, - ce qui n'est qu'une traduction du fait que les correspondances\inserted{cohomologiques} de Kunneth sur une variété algébrique projective lisse sont algébriques). On respire ! Mais on fera attention que les foncteurs $H^\bullet$ en question ne sont pas des foncteurs fibres proprement, mais des foncteurs fibres tordus.(On ne pourra s'empêcher de les appeler quand-même foncteurs fibres!) Quand on travaille avec des motifs pas nécessairement
